\section*{Research Foundations and Further Reading}

The attached predecessor essay, \emph{The Theory of Learning and Transformation} (25 October 2026), provides the primary conceptual foundation for this chapter, particularly its treatment of capability rather than content, prerequisite diagnosis, cognitive load, retrieval and spacing, guided practice, feedback, learner differences, assessment validity, AI assistance, and Bayes Institute's educational obligations. The present chapter reorganizes those ideas around \textbf{mechanisms of human transformation and the founder's decisions}, rather than reproducing the predecessor's extensive medium-by-medium instructional scripts. The following bibliography combines works developed in that essay with independently checked publisher, journal, institutional, and research records. A reference supports a \emph{bounded claim}; its inclusion should not be interpreted as an endorsement of every possible application of its theory.

\subsection*{Key Citations}

\begin{enumerate}[label={[\textbf{R\arabic*}]}]
    \item \textbf{National Academies of Sciences, Engineering, and Medicine (2018), \emph{How People Learn II: Learners, Contexts, and Cultures}.} Broad synthesis of learning, prior knowledge, motivation, metacognition, social environment, and lifelong development.
    \item \textbf{Kirschner, Sweller \& Clark (2006), ``Why Minimal Guidance During Instruction Does Not Work.''} Influential argument about cognitive architecture, working memory, and guidance for novices; useful within an ongoing debate about the design and timing of inquiry and discovery.
    \item \textbf{Kalyuga, Ayres, Chandler \& Sweller (2003), ``The Expertise Reversal Effect.''} Explains why instructional aids that help novices can become redundant or counterproductive as knowledge develops.
    \item \textbf{Chi \& Wylie (2014), ``The ICAP Framework: Linking Cognitive Engagement to Active Learning Outcomes.''} Theoretical and empirical account of passive, active, constructive, and interactive engagement; a framework with caveats rather than an instruction to maximize visible activity.
    \item \textbf{Roediger \& Karpicke (2006), ``Test-Enhanced Learning: Taking Memory Tests Improves Long-Term Retention.''} Experimental evidence that retrieval can improve delayed retention relative to additional restudy under tested conditions.
    \item \textbf{Dunlosky, Rawson, Marsh, Nathan \& Willingham (2013), ``Improving Students' Learning With Effective Learning Techniques.''} Comparative review of 10 common learning methods; especially influential for practice testing and distributed practice.
    \item \textbf{Cepeda, Pashler, Vul, Wixted \& Rohrer (2006), ``Distributed Practice in Verbal Recall Tasks.''} Review and quantitative synthesis showing that spacing effects depend on the relationship between learning intervals and retention horizons.
    \item \textbf{Pashler, McDaniel, Rohrer \& Bjork (2008/2009), ``Learning Styles: Concepts and Evidence.''} Review finding insufficient support for the common claim that matching fixed visual, auditory, or kinesthetic types to instructional modes improves learning.
    \item \textbf{Macnamara, Hambrick \& Oswald (2014), ``Deliberate Practice and Performance in Music, Games, Sports, Education, and Professions.''} Meta-analysis supporting a meaningful but variable role of practice and contradicting simplistic fixed-hour determinism.
    \item \textbf{Wisniewski, Zierer \& Hattie (2020), ``The Power of Feedback Revisited.''} Meta-analysis showing overall benefits with substantial variation and importance of information content.
    \item \textbf{Bloom (1984), ``The 2 Sigma Problem.''} Historically influential inquiry into tutoring, mastery learning, and whether group instruction can approach individualized instruction. Not a promise that modern tutoring universally produces 2 standard deviations of improvement.
    \item \textbf{Deci \& Ryan (2000), ``The 'What' and 'Why' of Goal Pursuits.''} Self-determination theory's treatment of autonomy, competence, and relatedness in motivation.
    \item \textbf{Deslauriers, McCarty, Miller, Callaghan \& Kestin (2019), ``Measuring Actual Learning Versus Feeling of Learning.''} A specific classroom experiment showing that subjective fluency and measured achievement can diverge.
    \item \textbf{Zimmerman (2002), ``Becoming a Self-Regulated Learner: An Overview.''} Theoretical framework for forethought, monitoring, strategic learning behavior, and reflection.
    \item \textbf{Barnett \& Ceci (2002), ``When and Where Do We Apply What We Learn? A Taxonomy for Far Transfer.''} Shows why transfer claims must specify how task and context change.
    \item \textbf{Bastani et al. (2025), ``Generative AI Without Guardrails Can Harm Learning: Evidence From High School Mathematics.''} Field experiment distinguishing performance aided by AI from later unassisted learning; outcome and safeguards depend on tutor design.
    \item \textbf{Kestin et al. (2025), ``AI Tutoring Outperforms In-Class Active Learning: An RCT Introducing a Novel Research-Based Design in an Authentic Educational Setting.''} A positive, context-specific randomized comparison with immediate outcomes in college physics.
\end{enumerate}

\subsection*{Additional Important Perspectives}

\begin{itemize}
    \item \textbf{Freeman et al. (2014), ``Active Learning Increases Student Performance in Science, Engineering, and Mathematics.''} Meta-analysis of undergraduate STEM instructional comparisons; a strong aggregate finding, not a guarantee for every activity, subject, learner, or instructor.
    \item \textbf{Wiggins \& McTighe (2005), \emph{Understanding by Design}.} Foundational argument for reasoning backward from desired performances and evidence.
    \item \textbf{Mislevy, Steinberg \& Almond (2003), ``On the Structure of Educational Assessments.''} Evidence-centered design linking competence claims, observable evidence, and assessment tasks.
    \item \textbf{Kane (2013), ``Validating the Interpretations and Uses of Test Scores.''} Argument-based validity: the meaning of a score depends on the strength of the inferences and proposed uses.
    \item \textbf{Theobald et al. (2020), ``Active Learning Narrows Achievement Gaps for Underrepresented Students in Undergraduate Science, Technology, Engineering, and Math.''} Review indicating that well-designed intensive active learning can improve some achievement-gap outcomes in the populations studied, not universally.
    \item \textbf{Mayer (2017), ``Using Multimedia for E-Learning.''} Research-informed design principles for reducing extraneous cognitive burden and improving learning from multimedia.
    \item \textbf{Noetel et al. (2021), ``Video Improves Learning in Higher Education.''} Systematic review of randomized comparisons, with important differences between substitution and supplementation and across designs.
    \item \textbf{CAST (2024), \emph{Universal Design for Learning Guidelines 3.0}.} A useful framework for considering access, representation, learner agency, and barriers, though specific implementations should be evaluated in context.
    \item \textbf{UNESCO, \emph{Guidance for Generative AI in Education and Research}.} Human-centered considerations including privacy, equity, oversight, and responsible adoption of generative AI.
\end{itemize}

\subsection*{Reading Research Without Turning It Into a Sales Claim}

These references differ in population, task, age, baseline knowledge, comparison condition, duration, environment, and outcome measure. Findings on immediate tests do not directly prove long-term professional transfer. Findings from school-age students do not automatically generalize to adult workers. A meta-analysis identifies average tendencies and heterogeneity; it does not determine the effect for a future Bayes cohort. A single impressive AI-tutoring trial does not settle the human-versus-AI debate. Self-determination theory is a framework for interpreting motivation, not a recipe that guarantees retention when 3 boxes are checked.

Treat research as a set of \textbf{well-informed priors, mechanisms to investigate, and cautions against naïve practice}. The institution still needs direct evidence that its own combination of teachers, tasks, learners, time, and resources creates the claimed capabilities.

\vspace{1.5em}
\noindent\emph{End of Chapter 11: The Theory of Learning and Human Transformation.}
